%% ma: L10-B14-0059
%% lop: 10
%% chuong: 5
%% bai: 14
%% dang: 10.14.8
%% loai: TN4
%% muc: VD
%% dap_an: "A"
%% nguon: "chuyen-de-cac-so-dac-trung-cua-mau-so-lieu-khong-ghep-nhom-toan-10-chuong-trinh-moi.pdf (GV. Phan Nhật Linh) (P05), Bài 03 Dạng 2 Phần I Câu 3"
%% kien_thuc: "Lớp 10 Bài 13 (số trung bình), Bài 14 (phương sai, độ lệch chuẩn)"
%% trang_thai: chinh-thuc
%% ghi_chu: "Tự duyệt theo yêu cầu giáo viên 22/10/2026."
\begin{ex}
Tiền thưởng (triệu đồng) cho cán bộ và nhân viên trong công ty X được trình bày trong bảng tần số sau:
\begin{center}
\begin{tabular}{|c|c|c|c|c|c|c|}\hline
Tiền thưởng $(x)$ & $2$ & $3$ & $4$ & $5$ & $6$ & \\\hline
Tần số $(n)$ & $5$ & $15$ & $10$ & $6$ & $7$ & $n=43$\\\hline
\end{tabular}
\end{center}
Phương sai của mẫu số liệu (làm tròn đến hàng phần trăm) là
\choice{\True $1{,}59$}{$1{,}58$}{$1{,}61$}{$1{,}57$}
\loigiai{$\overline x=\dfrac{167}{43}\approx3{,}884$; $s^2=\dfrac{1}{43}\sum n_ix_i^2-\overline x^2=\dfrac{717}{43}-\left(\dfrac{167}{43}\right)^2\approx1{,}59$. Chọn A.}
\end{ex}
